\SourceLeaf{058}%
porro. Hinc \(\Sigma((a\alpha+b\beta+c\gamma+\text{etc.})^2)\) necessario est minor quam \(\pi\). \textit{Secundo} habetur \(\Sigma(a\alpha+b\beta+c\gamma+\text{etc.})=\rho\), quoniam fit \(\Sigma a\alpha=1,\allowbreak\Sigma b\beta=1,\allowbreak\Sigma c\gamma=1\) etc.; unde facile deducitur, summam quadratorum \(\Sigma((a\alpha+b\beta+c\gamma+\text{etc.})^2)\) esse maiorem quam \(\frac{\rho\rho}{\pi}\), vel saltem non minorem. Hinc terminus
\[
\frac{\nu^4-3\mu^4}{(\pi-\rho)^2}.[\rho-\Sigma((a\alpha+b\beta+c\gamma+\text{etc.})^2)]
\]
necessario iacet inter limites \(-\frac{\nu^4-3\mu^4}{\pi-\rho}\) et \(\frac{\nu^4-3\mu^4}{\pi-\rho}.\frac{\rho}{\pi}\) vel, si latiores praeferimus, inter hos \(-\frac{\nu^4-3\mu^4}{\pi-\rho}\) et \(+\frac{\nu^4-3\mu^4}{\pi-\rho}\), et proin erroris medii in valore ipsius \(\mu\mu=\frac{M}{\pi-\rho}\) metuendi quadratum inter limites \(\frac{2\nu^4-4\mu^4}{\pi-\rho}\) et \(\frac{2\mu^4}{\pi-\rho}\), ita ut praecisionem quantamvis assequi liceat, si modo observationum multitudo fuerit satis magna.

Valde memorabile est, in hypothesi ea (art.\ 9, III), cui theoria quadratorum minimorum olim superstructa fuerat, illum terminum omnino excidere, et sicuti, ad eruendum valorem approximatum erroris medii observationum \(\mu\), in omnibus casibus aggregatum \(\lambda\lambda+\lambda'\lambda'+\lambda''\lambda''+\text{etc.}=M\) ita tractare oportet, ac si esset aggregatum \(\pi-\rho\) errorum fortuitorum, ita in illa hypothesi etiam praecisionem ipsam huius determinationis aequalem fieri ei, quam determinationi ex \(\pi-\rho\) erroribus veris tribuendam esse in art.\ 15 invenimus.
\par\nobreak\medskip\centerline{\rule{.25\linewidth}{.3pt}}
